\documentclass[12pt]{article}
\usepackage[utf8]{inputenc}
\usepackage[T1]{fontenc}
\usepackage{hyperref}
\title{Tanggu Language Resources for K-12 Curriculum Development}
\author{Compiled by Claude Sonnet 5 (claude-sonnet-5) for Andr\'as Kornai}
\date{}
\begin{document}
\maketitle

\section{Introduction}
Tanggu (ISO 639-3 tgu) is a Ramu-family language of the Bogia area, Madang Province, Papua New Guinea. Sources describe it as part of the small Ataitan (Moam River) family, with Igom as its closest relative. Speaker counts vary by source: about 3,000 (a 1991 figure cited on Wikipedia) and 2,330 (Endangered Languages). My searches found no school materials, primers or textbooks in Tanggu.

\section{Primary Educational Resources}
None were found in search results. No primers, textbooks or school programmes for Tanggu turned up.

\section{Linguistic and Reference Materials}
\begin{itemize}
\item Tanggu language, Wikipedia overview with classification, dialects and phonology: \url{https://en.wikipedia.org/wiki/Tanggu_language}
\item Ataitan languages, Wikipedia: \url{https://en.wikipedia.org/wiki/Ataitan_languages}
\item Tanggu page at the New Guinea World project, University of Zurich: \url{https://newguineaworld.linguistik.uzh.ch/families/keram-and-ramu-rivers/ramu-river/guam-and-moam-rivers/moam-river/igom-tanggu/tanggu}
\item Numeral systems entry for Tanggu, MPI-EVA: \url{https://lingweb.eva.mpg.de/channumerals/Tanggu.htm}
\end{itemize}

\section{Religious and Community Resources}
No scripture translations were found. Joshua Project has a people-group profile for the Tangu (it reports a population of 8,200, a different figure from the speaker counts above): \url{https://joshuaproject.net/people_groups/15254}

\section{Digital and Audio Resources}
No audio or video resources were found. The Endangered Languages Project has a catalogue entry that classes the language as threatened: \url{https://endangeredlanguages.com/lang/10653}

\section{Curriculum Development Context}
Little usable material exists. Any curriculum would have to begin with community-led documentation: a practical orthography, a wordlist and recorded texts. The sources above give only background and classification.
\end{document}